
%=====================================================
\section{Partially Ordered Sets}
%=====================================================

%-----------------------------------------------------

\begin{frame}{Partial Order}

\begin{block}{Definition}

\vspace{0.1cm}

A relation $R$ on a set $A$ is called a \textbf{Partial Order} if

\[
R
\text{ is }
\boxed{\text{Reflexive}}
\land
\boxed{\text{Antisymmetric}}
\land
\boxed{\text{Transitive}}
\]

\end{block}

\vspace{0.3cm}

A set together with a partial order,
\[
(A,R)
\]
is called a
\[
\boxed{\textbf{Partially Ordered Set (Poset)}}
\]

\end{frame}

%-----------------------------------------------------

\begin{frame}{Power Set as a Poset}

Let
\[
A=\{1,2,3\}
\]
Its power set is
\[
\mathcal P(A)
=
\{
\emptyset,
\{1\},
\{2\},
\{3\},
\{1,2\},
\{1,3\},
\{2,3\},
\{1,2,3\}
\}
\]
Define
\[
R=\subseteq
\]
Then
\[
(\mathcal P(A),\subseteq)
\]
is a Partially Ordered Set.

\end{frame}

%-----------------------------------------------------

\begin{frame}{Application}

A software project contains the tasks
\[
Collect
\rightarrow
Design
\rightarrow
Develop
\rightarrow
Test
\rightarrow
Deploy
\]

Define
\[
A\le B
\]
to mean
\[
\text{Task }A
\text{ must be completed before Task }B.
\]

This dependency relation is
\begin{itemize}
\item Reflexive
\item Antisymmetric
\item Transitive
\end{itemize}
Hence, it forms a Partial Order.

\end{frame}